\chapter{Waterfilling for Power Allocation}
\label{appendix:waterfilling}

In this appendix, we solve the problem of optimal power allocation accross independent scaler Gaussian channels using the method of Lagrangian multipliers.
The key principles of the method are explained alongside the derivation.
This method is also employed to solve other optimization problems encountered in this dissertation.
The solution to the optimal power allocation problem gives rise to the waterfilling algorithm, which is discussed at the end of this appendix.


\section{Optimal Power Allocation using Lagrangian Multipliers}
\label{app_section:optimal_power_allocation_using_lagrangian_multipliers}

Consider $N$ independent \ac{SISO} Gaussian trasmission channels.
Optimal power allocation across these channels can be formulated as a constrained maximization problem.
The objective is to distribute the total available transmit power among the different channels such that the combined achievable rate is maximized.

Consider the following general optimization problem:
\begin{equation}
\label{eq:optimal_power_allocation_optimization_problem}
\begin{aligned}
    & \underset{\{p_n\}}{\text{max}}
    & & \sum_{n=1}^{N} \, \alpha_n \log_2 \left( 1 + \gamma_n \, p_n \right) \\
    & \text{s. t.}
    & & \sum_{n=1}^{N} p_n = P_t \quad \text{and} \quad p_n \geq 0 \quad \forall n = 1, \dots, N
\end{aligned}
\end{equation}
where $p_n$ denotes the power allocated to channel $n$, and $P_t$ denotes the total available transmit power. 
The factors $\{\alpha_n\}$ represent weighting factors that allow for unequal prioritization of the individual transmission channels, while $\{gamma_n\}$ capture the ratio between the channel gain and the noise power of each indivudual channel.

To solve the constrained optimization problem in~\eqref{eq:optimal_power_allocation_optimization_problem}, the method of Lagrangian multipliers is employed.
The core idea of this method is to convert a constrained optimization problem into an unconstrained one by absorbing the constraints into a modified objective function, known as the Lagrangian.
This is achieved by introducing auxiliary variables, called Lagrange multipliers, one for each constraint, which penalize violations of the constraints.

Since the optimization problem in~\eqref{eq:optimal_power_allocation_optimization_problem} involves both an equality constraint (total power budget) and inequality constraints (non-negativity of the individual power allocations), both types must be incorporated into the Lagrangian.
Equality constraints enter with an unconstrained multiplier $\mu \in \mathbb{R}$, whereas inequality constraints of the form $p_n \geq 0$ are rewritten as $-p_n \leq 0$ and enter with a non-negative multiplier $\lambda_n \geq 0$, ensuring that the penalty is only active when the constraint is binding.
The resulting Lagrangian is
\begin{equation}
    \mathcal{L}(\mathbf{p}, \mu, \boldsymbol{\lambda}) = \sum_{n=1}^{N} \alpha_n \log_2 \left( 1 + \gamma_n \, p_n \right) - \mu \left(\sum_{n=1}^{N} p_n - P_t \right) - \sum_{n=1}^{N} \lambda_n p_n
\end{equation}
where $\mu$ is the Lagrange multiplier associated with the total power constraint, and $\lambda_n$ are the multipliers associated with the non-negativity constraints on the individual power allocations.

Since the objective function is concave and all constraints are affine, the optimization problem is convex.
For convex problems, the Karush-Kuhn-Tucker (KKT) conditions are both necessary and sufficient for a solution to be globally optimal~\cite{convex_optimization_2004, KKT_conditions_2009}.
These conditions are given by
\begin{subequations}
\label{eq:KKT_conditions}
\begin{align*}
    & \text{1. Stationarity:}               && \nabla_{\mathbf{p}} \, \mathcal{L}(\mathbf{p}, \mu, \boldsymbol{\lambda}) = 0 \\
    & \text{2. Primal feasibility:}         && p_n \geq 0  \quad \text{and} \quad \sum_{n=1}^{N} p_n = P_t \\
    & \text{3. Dual feasibility:}           && \lambda_n \geq 0 \\
    & \text{4. Complementary slackness:}    && \begin{cases} p_n > 0 \implies \lambda_n = 0 \\ p_n = 0 \implies \lambda_n \geq 0 \end{cases}
\end{align*}
\end{subequations}

From the KKT conditions, we can derive the optimal power allocation for each \ac{SISO} channel. The derivation proceeds by successively applying the stationarity, complementary slackness, and primal feasibility conditions.

The stationarity condition requires that the gradient of the Lagrangian with respect to each optimization variable vanishes at the optimum. This expresses the balance between the marginal gain in the achievable rate and the penalties introduced by the power constraints. We find:
\begin{equation}
\label{eq:KKT_stationarity_condition}
    \nabla_{\mathbf{p}} \, \mathcal{L}(\mathbf{p}, \mu, \boldsymbol{\lambda}) = 0
    \quad \iff \quad
    \frac{\partial \mathcal{L}}{\partial p_n} = \frac{\alpha_n \gamma_n}{\ln(2) \left( 1 + \gamma_n p_n \right)} - \mu - \lambda_n = 0, \quad \forall n = 1, \dots, N
\end{equation}

The complementary slackness condition links the non-negativity constraint on the allocated power to its associated Lagrange multiplier. When a positive amount of power is allocated to channel $n$, the corresponding multiplier $\lambda_n$ is zero, and the stationarity condition reduces to:
\begin{subequations}
\label{eq:KKT_complementary_slackness_condition}
\begin{align}
    & \frac{\alpha_n \gamma_n}{\ln(2) (1 + \gamma_n p_n)} - \mu = 0 \\
    & \iff \, p_n = \frac{\alpha_n}{\mu \ln(2)} - \frac{1}{\gamma_n}
\end{align}
\end{subequations}

The primal feasibility condition enforces the non-negativity of the power allocations. As a result, the above expression for $p_n$ must be clipped to zero whenever it becomes negative:
\begin{equation}
\label{eq:KKT_primal_feasibility_condition1}
    p_n = \left( \frac{\alpha_n}{\mu \ln(2)} - \frac{1}{\gamma_n} \right)^{+}
\end{equation}

Finally, the primal feasibility condition also requires the total allocated power to equal the available transmit power $P_t$. By substituting the expression for $p_n$ from Eq.~\eqref{eq:solution_pn} into this constraint, we obtain an implicit equation for the the Lagragian multiplier $\mu$:
\begin{equation}
\label{eq:KKT_primal_feasibility_condition2}
    \sum_{n=1}^{N} \left( \frac{\alpha_n}{\mu \ln(2)} - \frac{1}{\gamma_n} \right)^{+} = P_t
\end{equation}

Since $\mu$ cannot be obtained in closed form, it is determined iteratively such that the total allocated power satisfies \eqref{eq:KKT_primal_feasibility_condition2}. Owing to the monotonic relationship between $\mu$ and the total allocated power, a bisection method is typically employed to efficiently compute its optimal value.


\section{Waterfilling Algorithm}
\label{app_section:waterfilling_algorithm}

Applied to the more specific scenario considered in Section~\ref{section:capacity_achieving_precoding}, no prioritization is employed, i.e. $\alpha_n = 1$, and the \ac{SISO} Gaussian channels correspond to the eigen subchannels.

\begin{algorithm}[H]
\caption{Waterfilling Algorithm for power allocation across eigen subchannels}
\label{alg:waterfilling}
\begin{algorithmic}[1]
    \Require The total transmit power $P_t$; the SNR ratios $\gamma_i = \sigma_i / N_0$ for $i = 1, \dots, N_s$, sorted in decreasing order.
    \Ensure  The optimal power allocation $p_i$ for $i = 1, \dots, N_s$.
    
    \smallskip
    \Statex \hspace{-\algorithmicindent}\textbf{Step 1:} Determine the optimal number of active streams
    \State $N_s^{\mathrm{LB}}, \, N_s^{\mathrm{UB}} \leftarrow 0, \, N_s$
    \While{$N_s^{\mathrm{UB}} - N_s^{\mathrm{LB}} > 1$}
        \Statex \hspace{\algorithmicindent/2} Update the number of active streams: 
        \State \hspace{\algorithmicindent/2} $N_s^{\mathrm{iter}} \leftarrow \left\lceil (N_s^{\mathrm{UB}} + N_s^{\mathrm{LB}}) / 2 \right\rceil$
        \Statex \hspace{\algorithmicindent/2} Compute the minimum required power to activate $N_s^{\mathrm{iter}}$ streams:
        \State \hspace{\algorithmicindent/2} $P_{\mathrm{min}} \leftarrow \textstyle\sum_{i=1}^{N_s^{\mathrm{iter}}} \left( 1/\gamma_{N_s^{\mathrm{iter}}} - 1/\gamma_i \right)$
        \Statex \hspace{\algorithmicindent/2} Update the bound on the number of active streams:
        \State \hspace{\algorithmicindent/2} \textbf{if} $P_{\mathrm{min}} \leq P_t$ \textbf{then} $N_s^{\mathrm{LB}} \leftarrow N_s^{\mathrm{iter}}$ \textbf{else} $N_s^{\mathrm{UB}} \leftarrow N_s^{\mathrm{iter}}$
    \EndWhile
    
    \smallskip
    \Statex \hspace{-\algorithmicindent}\textbf{Step 2:} Compute the power allocated to each stream
    \Statex Fill the $N_s^{\mathrm{LB}}$ active stream containers:
    \State \indent $p_i \leftarrow 1/\gamma_{N_s^{\mathrm{LB}}} - 1/\gamma_i$ \quad for $i = 1, \dots, N_s^{\mathrm{LB}}$
    \Statex Distribute the remaining power equally among the active streams:
    \State \indent $p_i \leftarrow p_i + (P_t - \textstyle\sum_{i=1}^{N_s^{\mathrm{LB}}} p_i)\, /\, N_s^{\mathrm{LB}}$ \quad for $i = 1, \dots, N_s^{\mathrm{LB}}$
    \Statex Set the power of the inactive streams to zero:
    \State \indent $p_i \leftarrow 0$ \quad for $i = N_s^{\mathrm{LB}} + 1, \dots, N_s$
    
    \smallskip
    \Statex \hspace{-\algorithmicindent}\textbf{Return} $p_i$ for $i = 1, \dots, N_s$
\end{algorithmic}
\end{algorithm}